
\documentclass[12pt]{article}
\usepackage{amsmath, amsthm, amsfonts, mathtools, xcolor, caption}

\usepackage{minimalJ}

\numberwithin{equation}{section}

\usepackage[style=numeric,sorting=none,giveninits=true]{biblatex}
\addbibresource{bibliography.bib}

\usepackage{etoolbox}
\usepackage[margin=1in]{geometry}

\DeclareMathAlphabet{\mymathbb}{U}{BOONDOX-ds}{m}{n}

%\usepackage[colorlinks, allcolors=blue, hypertexnames=false]{hyperref}
%\usepackage[noabbrev, nameinlink]{cleveref}

\usepackage{comment}

%%%%%%%%%%%%%%%%%%%%%%%%%%%%%%%%%%%%%%%%%
%       Theorem Environments            %
%%%%%%%%%%%%%%%%%%%%%%%%%%%%%%%%%%%%%%%%%

\newtheorem{theorem}{Theorem}[section]
\newtheorem{lemma}[theorem]{Lemma}
\newtheorem{corollary}[theorem]{Corollary}
\newtheorem{proposition}[theorem]{Proposition}
\theoremstyle{definition}
\newtheorem{remark}{Remark}
\newtheorem{definition}{Definition}
\newtheorem{notation}{Notation}
\newtheorem{example}{Example}
\newtheorem{conjecture}[theorem]{Conjecture}
\newtheorem{openproblem}[theorem]{Open Problem}

%%%%%%%%%%%%%%%%%%%%%%%%%%%%%%%%%%%%%%%%%
%               New Commands            %
%%%%%%%%%%%%%%%%%%%%%%%%%%%%%%%%%%%%%%%%%

\newcommand{\ev}[1]{( #1 )} 
\newcommand{\ew}[1]{\left( #1 \right)} 

\newcommand{\N}{\mathbb{N}}
\newcommand{\Z}{\mathbb{Z}}
\newcommand{\C}{\mathbb{C}}
\newcommand{\Q}{\mathbb{Q}}
\newcommand{\R}{\mathbb{R}}



%%%%%%%%%%%%%%%%%%%%%%%%%%%%%%%%%%%%%%%%%
%               Document                %
%%%%%%%%%%%%%%%%%%%%%%%%%%%%%%%%%%%%%%%%%

\let\underbrace\LaTeXunderbrace
\let\overbrace\LaTeXoverbrace


\begin{document}

\section*{Challenge 8: Ryser's Hypergraph Conjecture}
\refstepcounter{section}



\begin{definition}[$r$-partite]
Given $r\in \mathbb{N}$, an $r$-uniform hypergraph $\mathcal{H}$ is \emph{$r$-partite} if there is a partition
$$
V\ev{\mathcal{H}} = \bigcup_{i=1}^{r} P_i
$$
of the vertex set of $\mathcal{H}$ such that $|P_i \cap e| = 1$, for each $e\in E\ev{\mathcal{H}}$ and $i\in [r]$. 
\end{definition}


\begin{definition}[Cover number]
Given a hypergraph $\mathcal{H}$, a
set $C$ of vertices of $\mathcal{H}$ is a \textit{vertex cover} if it contains a vertex of every edge. The smallest cardinality of a vertex cover of $\mathcal{H}$ is called the \textit{cover number} of $\mathcal{H}$. It is denoted by $\tau(\mathcal{H})$.
\end{definition}

\begin{definition}[Matching number]
The \textit{matching number} of a hypergraph $\mathcal{H}$ is the largest cardinality of a set of pairwise disjoint edges of $\mathcal{H}$. It is denoted by $\nu\ev{\mathcal{H}}$.
\end{definition}

\begin{conjecture}[Ryser's Conjecture \mbox{\cite{bishnoi}}]
Let $r\in \mathbb{N}$ with $r\geq 2$.
Every $r$-partite $r$-uniform hypergraph $\mathcal{H}$ satisfies:
$$
\tau\ev{\mathcal{H}} \leq \ev{r-1} \cdot \nu\ev{\mathcal{H}}
$$

\end{conjecture}

For $r=2$, it is a classical theorem of K\H{o}nig, and for $r=3$ it was proved by Aharoni \cite{AharoniRyserTripartite} using the topological method. The hypothesis $r\geq 2$ is necessary: for $r=1$ the bound reads $\tau\ev{\mathcal{H}}\leq 0$, which already fails for the single-vertex single-edge hypergraph (where $\tau = \nu = 1$).


%%%%%%%%%%%%%%%%%%%%%%%%%%%%%%%%%%%%%%%%%%%%%%%%%%
\printbibliography
%%%%%%%%%%%%%%%%%%%%%%%%%%%%%%%%%%%%%%%%%%%%%%%%%%
\end{document}
